\subsection{محدودیت‌ها و تحلیل شکست}
بررسی عملکرد ماژول \lr{aegis\_v6} چالش‌های ساختاری زیر را تبیین می‌نماید:
\begin{itemize}
    \item \textbf{علت قرارگیری در رده \lr{PASS\_LOW} در مدل‌های انکودر و دکودر:} دقت \lr{52.87\%} در \lr{DistilBERT}، \lr{50.92\%} در \lr{RoBERTa} و \lr{50.46\%} در \lr{GPT-2} نشان‌دهنده پدیده انجماد یادگیری لایه‌های توجه خطی در حضور نگاشت کِرنل متعامد است. در فاز انطباق وزن‌ها با روش بازآموزی با گام‌های محدود، کِرنل غیرخطی تصادفی به دلیل بالا بودن واریانس پایدارسازهای نمایی و توزیع مقادیر ویژه، ظرفیت تفکیک الگوهای معنایی ظریف در وظیفه تحلیل احساسات را به شدت تحدید کرده و رفتار مدل را به پیش‌بینی اکثریت کلاس‌ها میل می‌دهد.
    \item \textbf{تحلیل عملکرد در مدل \lr{Qwen} و خطاهای سازگاری:} در مدل زبانی \lr{Qwen} با دقت \lr{47.71\%} وضعیت قبولی \lr{PASS} احراز شده که نشان از حفظ پایداری محاسبات خودبازگشتی دارد، اما به دلیل عدم انطباق لایه‌های فاکتورگیری با فرمت ماسک‌های سه بعدی در مدل \lr{T5}، ارزیابی روی آن با خطای عدم انطباق (\lr{SKIP\_T5\_INCOMPATIBLE}) مواجه شده است.
\end{itemize}
